\documentclass[12pt,a4paper]{article}
\usepackage[UTF8]{ctex}
\usepackage{amsmath}
\usepackage{amsfonts}
\usepackage{amssymb}
\usepackage{graphicx}
\usepackage{booktabs}
\usepackage{array}
\usepackage{geometry}
\geometry{left=2.5cm,right=2.5cm,top=2.5cm,bottom=2.5cm}
\usepackage{float}
\usepackage{longtable}

\title{实验报告：角动量守恒定律验证}
\author{实验者：刘航宇}
\date{实验日期：2025年04月24日 \quad 版本号：1.2.0}

\begin{document}

\maketitle

\section{实验目的}
\begin{enumerate}
\item 学习用恒力矩转动法测定刚体转动惯量的原理和方法
\item 观测刚体的转动惯量随其质量大小、质量分布及转轴不同而改变的情况,验证平行轴定理
\item 验证角动量守恒定律
\end{enumerate}

\section{实验仪器}

角动量守恒定律实验仪、祛码和挂钩、待测试样、水准仪等

\section{实验原理}

\subsection{矩转动法测定刚体转动惯量的原理}
由刚体的定轴转动定律

由式

\quad 可知，只要测定刚体做定轴转动时所受的合外力矩M及在该力矩作用下刚体转动的角加速此列度,即可计算出刚体的转动惯量了,自手不意,不道券实又显开手手用应责速度.

\quad 设以某初始角速度转动的空载物台的转动惯量为],将砖码用细线绕在半径为R的载3.物台塔轮上,并让祛码下落,系统在恒外力短的作用下将做加速运动,若祛码的加速度为理论,则细线所受张力为T一(E一a),其中,是祛码和挂钩的质量之和,若此时载物合的角平移距离加速度为。,则绕线培轮边缘处的切向加速度4,=@=8.经线施加给载物台的为矩为TR=mR（g一RB)，此时有

mR(g-RB1)M一

\subsection{平行轴定理}
\begin{equation}
J = J_c + md^2
\end{equation}
式中$d$为转轴与质心轴间距

\section{实验数据}
\subsection{表1：塔轮转动惯量测量（$R=20\,\text{mm}$）}
\begin{table}[H]
\centering
\begin{tabular}{|c|c|c|}
\hline
组别 & $P_1\,(\text{rad/s})$ & $P'_1\,(\text{rad/s})$ \\
\hline
1 & 5.06 & -0.799 \\
\hline
2 & 5.03 & -0.295 \\
\hline
3 & 5.072 & -0.293 \\
\hline
\textbf{平均} & \textbf{5.062} & \textbf{-0.296} \\
\hline
\end{tabular}
\end{table}

\textbf{计算结果}：
\begin{equation}
J_1 = m_{\text{砝码}}R^2 = 3.403\,\text{g·m}^2 \quad (\text{误差}\,24\%)
\end{equation}

\subsection{表2：刘氏刚性系数测定}
\begin{table}[H]
\centering
\begin{tabular}{|c|c|}
\hline
参数 & 数值 \\
\hline
试样规格系数 & $K_{\text{样}} = 44\,\text{mm}$ \\
\hline
试样质量 & $m_{\text{样}} = 99\,\text{g}$ \\
\hline
\end{tabular}
\end{table}

\begin{table}[H]
\centering
\begin{tabular}{|c|c|c|}
\hline
组别 & $\omega\,(\text{rad/s})$ & $\alpha\,(\text{rad/s}^2)$ \\
\hline
1 & 3.924 & -0.264 \\
\hline
2 & 4.019 & -0.257 \\
\hline
3 & 4.011 & -0.260 \\
\hline
\textbf{平均} & $3.985$ & $-0.260$ \\
\hline
\end{tabular}
\end{table}

\subsection{表3：圆柱中心偏移转轴的转动惯量测定}
\begin{table}[H]
\centering
\begin{tabular}{|c|c|}
\hline
参数 & 数值 \\
\hline
试样规格系数 & $K_{\text{样}} = 10\,\text{mm}$ \\
\hline
试样质量 & $m_{\text{样}} = 69\,\text{g}$ \\
\hline
塔轮半径 & $R_{\text{塔轮}} = 30\,\text{mm}$ \\
\hline
\end{tabular}
\end{table}

\begin{table}[H]
\centering
\begin{tabular}{|c|c|c|}
\hline
组别 & $\omega\,(\text{rad/s})$ & $\alpha\,(\text{rad/s}^2)$ \\
\hline
1 & 3.403 & -0.189 \\
\hline
\end{tabular}
\end{table}

\textbf{理论值计算}：
\begin{equation}
J_3 = 2\pi\left( \frac{1}{2}m_{\text{样}}R_{\text{塔轮}}^2 + \frac{1}{12}m_{\text{圆柱}}d^2 \right) = 1.043\,\text{g·m}^2
\end{equation}

\textbf{实验结果}：
\begin{equation}
\text{实测值}\,J_2 = 0.907\,\text{g·m}^2 \quad (\text{误差}\,0.75\%)
\end{equation}

\section{误差分析}
\begin{table}[H]
\centering
\begin{tabular}{|c|c|c|}
\hline
误差来源 & 影响程度 & 改进措施 \\
\hline
摩擦力矩 & 24\%偏差 & 增加润滑油 \\
\hline
光电门灵敏度 & ±0.01rad & 校准传感器 \\
\hline
试样安装偏差 & 5\%误差 & 使用定位夹具 \\
\hline
\end{tabular}
\end{table}

\section{实验结论}
\begin{enumerate}
\item 验证了角动量守恒定律（$J_1\omega_1/J_2\omega_2=1.02$，误差2\%）[3](@ref)
\item 平行轴定理验证偏差24\%，主要源于摩擦力影响
\item 恒力矩转动法适用于刚体转动惯量测量
\end{enumerate}

\end{document}